\sectiontitle{参考答案}

\textbf{课后练习1（A组）}
\begin{enumerate}
\item (a) 普查；(b) 抽样调查；(c) 抽样调查；(d) 抽样调查
\item 总体：$1500$名学生；个体：每名学生；样本：$100$名学生
\item $m=60-(10+15+20)=15$，频率$=15/60=0.25$
\item 优秀$6/40=0.15$，良好$14/40=0.35$，合格$16/40=0.4$，不合格$4/40=0.1$
\item 总数$=12\div0.25=48$
\item 频率$=18/50=0.36$，圆心角$=360^\circ\times0.36=129.6^\circ$
\end{enumerate}

\textbf{课后练习1（B组）}
\begin{enumerate}
\item 方案示例：按区分层抽样，抽取$500$户。普查需查$50000$户，成本太高且没必要
\item 红：$50\times(8/20)=20$；黄：$50\times(7/20)=17.5\approx18$；蓝：$50\times(5/20)=12.5\approx13$
\item 方案二（匿名）更可靠。方案一中，学生可能因为不好意思而不报真实的短睡眠时间
\item 回答问卷的$24\%$本身不是随机样本——愿意回答的人和不愿意回答的人可能有系统性的差异（比如对选举特别热情的人更可能回信）
\item 诱导版："为了校园安全，你是否支持禁止学生带手机？"中立版："你对中学生带手机入校持什么态度？（支持/反对/中立）"
\end{enumerate}

\textbf{课后练习2（A组）}
\begin{enumerate}
\item 春季$10/40=25\%$，$90^\circ$；夏季$12/40=30\%$，$108^\circ$；秋季$8/40=20\%$，$72^\circ$；冬季$10/40=25\%$，$90^\circ$
\item 周六（80）最好，周四（150）最差
\item 条形图直条间有空隙（类别独立），折线图点间有线连接（顺序关系）
\item 条形图（比较各季度）或折线图（看季度趋势）
\item $360^\circ$，因为各部分百分比之和为$100\%$
\item 近视$30/50=60\%$，$216^\circ$；非近视$20/50=40\%$，$144^\circ$
\end{enumerate}

\textbf{课后练习2（B组）}
\begin{enumerate}
\item 原A圆心角$=360^\circ\times30/60=180^\circ$；加D=20后总数$=80$，A圆心角$=360^\circ\times30/80=135^\circ$（A占比变小）
\item (a) 逐季上升；(b) Q2: $20\%$，Q3: $25\%$，Q4: $20\%$；(c) Q2到Q3最陡，增长率最大
\item (a) 零食$100$元，文具$50$元，交通$30$元，储蓄$20$元；(b) 零食$30\%$则储蓄$=100\%-30\%-25\%-15\%=30\%$
\item (a) 折线图；(b) 条形图；(c) 扇形图
\item 纵轴从$0$开始，$80$到$100$的差异在图中只占一小段，视觉上平缓很多。纵轴截断会夸大变化幅度
\end{enumerate}

\textbf{课后练习3（A组）}
\begin{enumerate}
\item (a) 平均$=9$，中位$=9$，无众数；(b) 平均$\approx5.83$，中位$=5.5$，众$=5$；(c) 全$=10$
\item 平均$=83$，中位$=(85+88)/2=86.5$
\item 排序$3,7,8,10,x$，中位$=8$则$x\ge8$，最小$=8$
\item A组：平均$=5$，中位$=5$，众$=5$；B组：平均$=5$，中位$=5$，无众数
\item 所有数据相等
\item $88\times0.4+76\times0.2+92\times0.4=35.2+15.2+36.8=87.2$分
\end{enumerate}

\textbf{课后练习3（B组）}
\begin{enumerate}
\item 原总身高$=29\times155=4495$，加入后$(4495+175)/30\approx155.67$cm
\item $a+b+c+d=40$，$b+c=24$，$a+d=16$，满足$a<b<c<d$正整数即可（如$3,11,13,13$）
\item (a) 平均$=(120+80+100)/20=15$万，中位$=8$万；(b) 公司选平均数$15$万宣传，求职者应看中位数$8$万
\item 示例：$10,13,14,14,9$（平均$=12$，中位$=13$，众$=14$）
\item 两次低分是极端值，把平均分拉低了。可以改用中位数，或者去掉最高最低后再求平均
\end{enumerate}

\textbf{课后练习4（A组）}
\begin{enumerate}
\item $\bar x=10$，$s^2=18$，$s=\sqrt{18}\approx4.24$
\item $s=3$；加常数后方差不变仍为$9$
\item (a) A: $\bar x=8$，$s^2=18$；B: $\bar x=8$，$s^2=2$；(b) B组更集中
\item 甲命中率$0.7$，方差$0.21$；乙命中率$0.6$，方差$0.24$；甲更稳定
\item $8$（数据乘$k$，标准差变为$|k|$倍）
\end{enumerate}

\textbf{课后练习4（B组）}
\begin{enumerate}
\item 加$a$：每个数据与平均数的离差不变，方差不变。乘$k$：每个离差乘$k$，方差乘$k^2$
\item 例：A组$1,2$（$\bar x=1.5$，$s^2=0.25$），B组$9,10$（$\bar y=9.5$，$s^2=0.25$），合并$1,2,9,10$，方差$=12.5\neq0.25$
\item 去掉极值后方差变小。体育比赛用此法避免评委个人偏见影响
\item $y_i=x_i-b$的方差仍为$s^2$；$z_i=kx_i$的方差为$k^2s^2$
\item 乙虽然平均命中率略低，但波动更小——关键时刻更可能发挥正常水平，不会突然失常
\end{enumerate}

\textbf{课后练习5（A组）}
\begin{enumerate}
\item (a) 随机事件；(b) 不可能事件；(c) 随机事件；(d) 必然事件
\item $5/8$
\item 既是奇数又是质数：$3,5$（$1$不是质数），$P=2/5$
\item $1/4$
\item $2/6=1/3$
\item $42/50=0.84$
\end{enumerate}

\textbf{课后练习5（B组）}
\begin{enumerate}
\item 四种：男男、男女、女男、女女；性别不同$P=1/2$
\item $C_4^2=6$种，$P=1/6$
\item 甲掷出$5$点：$P=1/6$；甲大于乙：列表得$15$种，$P=15/36=5/12$
\item 小红$0.495$更接近$0.5$，说明次数越多频率越接近概率
\item 数字：$1/10^4=0.0001$；字母：$1/26^4\approx0.0000022$
\item 公平方案：和为$7$（$6$种）甲赢，和为$6$或$8$（各$5$种共$10$种）乙赢——但$6\neq10$。更公平：和为$7$甲赢（$6$种），和为$2$或$12$乙赢（$2$种），其余重掷——此方案甲$P=1/6$，乙$P=1/18$，仍不公平。完全公平的方案需双方各$18$种，如甲：$2,3,4,10,11,12$（共$1+2+3+3+2+1=12$种），乙：$5,6,7,8,9$（$4+5+6+5+4=24$种）——也不等。真正的公平需调整计分而非单纯看胜负
\item 都正$P=1/4$，小明期望$=3/4=0.75$；都反$P=1/4$，小明期望$=1/4=0.25$；一正一反$P=1/2$，小红期望$=1$。小红期望更高，游戏不公平
\end{enumerate}

\textbf{课后练习6（A组）}
\begin{enumerate}
\item $10000$人中患者$10$人，真阳性$\approx10$，健康人$9990$人，假阳性$\approx100$，阳性共$\approx110$，患病概率$\approx10/110\approx9.1\%$
\item 回收率未知——可能只有满意的人填了，不满意的人没填
\item 不能，两者正相关可能受共同因素（医疗进步、生活水平）影响
\item 纵轴不从$0$开始会夸大差异
\item 根据数据：晴/高/非节假日——已有的$2$条晴/高数据都去了，且晴/非节假日也去了，预测"去"
\end{enumerate}

\textbf{课后练习6（B组）}
\begin{enumerate}
\item 方案B更好（分层抽样），保持男女比例与总体一致。A是简单随机抽样
\item $98\%$不可信，仅$5\%$买家评价（样本偏差）。应强制评价或按购买次数加权
\item 都正$P=1/4$，小明期望$=3/4=0.75$；都反$P=1/4$，小明期望$=1/4=0.25$；一正一反$P=1/2$，小红期望$=1$。游戏不公平
\item 略（开放性项目）
\item 略（开放性）
\end{enumerate}
